\documentclass[10pt]{article}
\usepackage[margin=20mm]{geometry}
\usepackage{amsmath,amssymb,mathtools,fontspec,unicode-math,xurl,hyperref}
\IfFontExistsTF{DejaVu Serif}{\setmainfont{DejaVu Serif}}{\setmainfont{Latin Modern Roman}}
\setmathfont{latinmodern-math.otf}
\setlength{\parindent}{0pt}\setlength{\parskip}{6pt}\sloppy
\title{Technical Verification of
Matroids, Feynman Categories,
and Koszul Duality}
\author{AI-assisted technical review}
\date{26 October 2026}
\begin{document}
\maketitle
Abstract: Small geometric lattices exhibit failures in the printed normal-form invariant, normal-frame parametrization, cubical reordering formula, projective Leray differential, and the all-object projective Orlik–Solomon assignment. These are demonstrated problems in core constructions; explicit repairs are supplied with their precise limits. Five additional local issue groups concern an irreducibility restriction, a weak inequality, an unnecessary fiber equality, outlook axioms, and a generator index. A shifted-sign calculation is kept separate because its conclusion depends on the target symmetry. None of the examples disproves the intended Koszulness theorem. The complete external Leray comparison and global suspension compatibilities are not independently certified. Scientific decision: pending.\par
\section*{Exact version, mathematical scope, and evidence}
The source is Matroids, Feynman categories, and Koszul duality, by Basile Coron [paper].\par
Reviewed version: arXiv:2211.12370v3, dated 11 December 2024; 70 pages, 570554 bytes. Assigned journal DOI:\par
https://doi.org/10.1215/00127094-2024-0077.\par
The exact supplied author manuscript was reviewed; the journal version and later corrections were not compared. Source locations below are one-based physical PDF pages, agreeing with inspected printed labels. The source SHA256 is\par
`aa26c729c980c32c019fed5761f4367e81eb1a0e7635ede59316746481592d42`.\par
The main objects are finite geometric lattices with building sets. An irreducible built lattice means that \(\hat 1\in\mathcal G\); its underlying lattice need not be an indecomposable lattice. In particular, the Boolean examples below with top-containing building sets are permitted. The paper constructs a Feynman category, FY and OS (co)operadic structures, a quadratic presentation and Gröbner basis for \(\mathrm{FY}^\vee\), and two Koszul arguments: one through a Leray model, another through shuffle operads.\par
A native helper read all 70 pages, including the bibliography, and inspected 23 source images. Two other native helpers provided disjoint local cross-checks of the normal-form proof and the Leray correction. The lead reviewer read decisive passages on 27 pages, inspected 18 source images, reconstructed the examples and repairs, and ran an independent exact calculation. These roles are recorded separately; they are not human reviewers or external controllers. Full manifest parsing, hashing and rendering were not counted as scientific reading or image inspection. No author code was run.\par
The Major labels below concern invalid core derivations or definitions even when a short repair exists. Minor labels concern local statements whose needed replacements are available without changing the main mechanism. The sign concern in Section 8 (\#sec:sign) is qualified separately. Author adoption of every proposed change is unknown.\par
\section*{Major F1: the normal-form invariant does not descend}
\subsection*{A permitted composition with two invariant values}
Lemma 3.20 (p. 27) assigns raw inclusion to structural arrows on all non-bottom interval elements, inverse underlying poset maps to isomorphisms, and then extends by composition. Checking relations (6) and (7) does not check the additional structural composition (3) with its source identification \(\Phi\) from Lemma 3.12.\par
Let \(L=B_3\) have atoms \(a,b,c\), write \(ab=a\vee b\), and put \(t=abc\). Take \[\mathcal G=\{a,b,c,t\},\qquad S=\{a,t\}.\] For a proper Boolean element, its maximal building factors are its atoms and their product is its lower interval. For \(t\), the unique maximal factor is \(t\) itself. Thus \(\mathcal G\) is a building set, and \((L,\mathcal G)\) is irreducible. The upper interval \([a,t]\) has induced building set \(\{ab,ac,t\}\). Refine it by \(T=\{ab,t\}\), and use the unit in \([0,a]\). Both local sets are spanning and nested. Since \(\operatorname{Comp}_a(ab)=b\), equation (3) gives the composite \(U=\{a,b,t\}\), which is nested because \(ab\notin\mathcal G\).\par
The iterated and direct sources are respectively \[W=[0,a]\otimes[a,ab]\otimes[ab,t],\qquad
 W'=[0,a]\otimes[0,b]\otimes[ab,t].\] On the middle factor the underlying poset map defining \(\Phi\) is \(j:[0,b]\to[a,ab]\), \(j(x)=a\vee x\). The categorical isomorphism is represented in the opposite direction, so the proposed invariant sends \(ab\) to \(b\) along \(\Phi\). The two structural inclusions instead leave \(ab\) unchanged. Thus two expressions for the same arrow give \[\iota(S)\iota(\mathrm{Id}\otimes T)(ab)=ab,\qquad
 \iota(U)\iota(\Phi)(ab)=b.\] The invariant is not well-defined. In the maximal building set, \(\operatorname{Comp}\) happens to be inclusion; that special case cannot justify the general statement.\par
\subsection*{Complete replacement using building sets}
Equation reference: eq:compmap. Associate to an object its building set, with tagged disjoint union for tensor products. For a structural arrow \(S\), use \[
 c_S:\coprod_{G\in S}\operatorname{Ind}_{[\tau_S(G),G]}(\mathcal G)\longrightarrow\mathcal G,
 \qquad Y\longmapsto\operatorname{Comp}_{\tau_S(G)}(Y).\] For an isomorphism use the inverse poset map restricted to building sets. The domain restriction matters: \(\operatorname{Comp}\) need not be defined on every non-bottom lattice element.\par
Equation reference: eq:compidentity. Here is the full composition check. For one block of a refinement put \[A=\tau_S(G_i),\quad H=[A,G_i],\quad
 B=\tau_{T_i}(K),\quad C=\operatorname{Comp}_A^L(K),\quad D=\tau_U(C).\] The interval identification is \(j:[D,C]\to[B,K]\), \(j(Z)=A\vee Z\). The pointwise identity proved in Lemma 3.13 (pp. 22–24) is \[
 \operatorname{Comp}_A^L\!\left(\operatorname{Comp}_B^H(Y)\right)
   =\operatorname{Comp}_D^L\!\left(j^{-1}(Y)\right)\] for every \(Y\) in the induced building set of \([B,K]\). It follows by joining the proposed value with \(D\) and using the maximality characterization of \(\operatorname{Comp}\); the building-factor product decomposition proves that every other candidate lies below it. This preliminary proof precedes the normal-form theorem. Equation [eq:compidentity] (\#eq:compidentity) is exactly the missing structural compatibility.\par
The remaining relations hold as follows. A built-lattice isomorphism preserves the maximal defining fiber of \(\operatorname{Comp}\), so \(h\operatorname{Comp}_A(Y)=\operatorname{Comp}_{h(A)}h(Y)\), proving naturality. Permuting source factors permutes tagged blocks in [eq:compmap] (\#eq:compmap). Units use \(\operatorname{Comp}_0(Y)=Y\). Disjoint union gives tensor interchange and symmetry. Inverse poset maps respect identity and composition. Together these checks define \(c\) on the category as presented; they do not assume the later generator-presentation theorem.\par
For uniqueness of normal forms, labelled source-factor tops map to the corresponding nodes of \(S\). They recover \(S\) and the factor assignment; the chosen order determines its permutation. If \(h_i:L_i\to[\tau_S(H_i),H_i]\) is the forward interval isomorphism, recover it on building elements by \[h_i(g)=\tau_S(H_i)\vee c_F(g).\] Every atom lies in the building set, and an atomic lattice is join-generated by its atoms. Thus \(h_i\), including its bottom, is determined everywhere. Simply restricting \(c_F\) without joining the lower endpoint would not recover \(h_i\). Existence follows by moving isomorphisms to the source, flattening structural composites, and sorting the factors. This proves the intended normal form without circularity. It repairs the proof used for the presentation and shuffle comparison; it does not demonstrate a false normal-form theorem.\par
\section*{Major F2: a rank-one top interval is missing}
Proposition 5.23 and Lemma 5.24 (p. 53) prohibit every rank-one interval in the spanning frame. Yet their proposed frame is \[\operatorname{fr}(S)
 =\{G\in S:\operatorname{rk}[\tau_S(G),G]>1\}\cup\{t\}.\] The appended top may have local rank one.\par
Already a rank-one irreducible arity has its nonzero FY class and trivial normal monomial, but has no allowed frame under the printed restriction. A positive-weight example is \(B_3\) with maximal building set and atom order \(a<b<c\). The nested set \(\{ab,t\}\) is normal: its lower node has rank two, whereas the quadratic leading terms on p. 54 have a nonminimal atom as lower node. Its top interval \([ab,t]\) has rank one. Rank additivity implies that a rank-three frame with every interval of rank at least two can have only its top node, so the printed parametrization misses this monomial.\par
The correction is to prohibit rank-one intervals only at non-top frame nodes. Keep the displayed frame formula. Removed rank-one nodes form clusters in the frame intervals; each is assigned to its unique nearest retained parent. A proper retained node has residual rank greater than one, so its cluster is proper. The top cluster may exhaust all but one unit of its interval rank. Accordingly the existing bounds \[k_G<\operatorname{rk}[\tau(G),G]-1\quad(G\ne t),\qquad
 k_t\leq\operatorname{rk}[\tau(t),t]-1\] remain correct; a rank-one top interval allows \(k_t=0\). The exponent correspondence \[\alpha_G=\operatorname{rk}[\tau(G),G]-k_G-1\] then gives positive proper-node exponents and a possibly zero top exponent, exactly the spanning version of the FY basis. The intrinsic frame criterion and unique local clusters give uniqueness. With the minimum argument in Section 7 (\#sec:minor), the normal-cluster/EL-chain proof survives.\par
Exact enumeration in \(B_3\) gives six FY normal monomials, with half-degree Hilbert polynomial \(1+4z+z^2\). The printed frames give three; the corrected frames give six. Thus the printed basis-count step is false, but the corrected step supports the intended Gröbner conclusion.\par
\section*{Major F3: cubical insertion must use the current interval}
The formula on p. 57 joins the next node with earlier selected nodes strictly below it. That join does not change the node. For \(B_3\), \(\mathcal G=\{a,b,c,t\}\), and \(S=\{a,b,t\}\), insert \(a\) first. The current factors are \([0,a]\) and \([a,t]\). The printed next node \(b\) lies in neither factor. The needed induced node is \(ab\in[a,t]\), with \(\operatorname{Comp}_a(ab)=b\). The tensor generator asserted by the printed formula is therefore undefined.\par
Equation reference: eq:insert. A complete insertion rule is available. Start with \(T=\{t\}\). For each next original node \(G\), let \(P\) be the least element of \(T\) strictly above \(G\), and put \[
 A=\tau_T(P),\qquad H=A\vee G.\] Insert \(H\) in the factor \([A,P]\), then adjoin \(G\) to \(T\). The nested-factor decomposition gives \(A<H<P\), \(H\in\operatorname{Ind}_{[A,P]}(\mathcal G)\), and \(\operatorname{Comp}_A(H)=G\). Indeed \(G\leq A\) would put it below a smaller current parent; \(A\vee G=P\) would make a forbidden nested antichain with join in \(\mathcal G\). The factor characterization supplies the \(\operatorname{Comp}\) identity.\par
After the normal-form repair, chains modulo intermediate isomorphisms are precisely orderings of the non-top nodes. The rule [eq:insert] (\#eq:insert) reconstructs every ordering. Permuting those orderings gives a free symmetric-group action, preserves the composite, and has exactly the degree-\(n\) composition fibers as its orbits. Concatenation gives the required compatibility. The target in Definition 6.1 should accordingly be the degree-\(n\) component \(\operatorname{Hom}_n(A,B)\). The least-atom convention also respects these interval identifications in the shuffle category. The intended cubicality is supported by this replacement; the old formula does not prove it.\par
\section*{Major F4: the printed Leray differential is not defined}
\subsection*{A rank-two witness}
Definition 6.4 (pp. 61–62) includes all \(e_G,x_G\), kills \(e_t\), and prescribes \(d(e_G)=x_G\) without a top exception. A differential on a quotient must preserve its defining ideal.\par
Use the flats of the simple connected matroid \(U_{2,3}\): bottom, atoms \(a,b,c\), and top \(t\), with maximal building set. A nested set contains at most one atom. The atom-linear relations give \(x_a=x_b=x_c=-h\), where \(h=x_t\). The nonnested relations give \(h^2=0\), \(e_i h=0\), and \(e_i e_j=0\) for distinct atoms. In the graded quotient by \(e_t\), \(h\) is nonzero: the algebra map \[x_t\mapsto h,\qquad x_a,x_b,x_c\mapsto-h,\qquad e_G\mapsto0
 \quad\text{to }\mathbb Q[h]/(h^2)\] respects every defining relation. But the prescribed differential sends the zero class \(e_t\) to this nonzero class. Thus the object used immediately for the comparison of complexes is not a DGA as printed.\par
\subsection*{A well-defined correction and its projective map}
Equation reference: eq:bpr. Define instead \[
 B_{\mathrm{pr}}
 =\left(\Lambda(e_G:G\in\mathcal G\setminus\{t\})
          \otimes\mathbb Q[x_G:G\in\mathcal G]\right)/J,\] where \(J\) has the same nonnested mixed-support relations and atom-linear relations. Keep every \(x_G\), including \(x_t\). Set \(de_G=x_G\) for proper \(G\) and \(dx_G=0\). The bidegrees remain \((0,1)\) and \((2,0)\); \(d\) has bidegree \((2,-1)\).\par
Differentiating a nonnested monomial replaces one \(e_G\) by \(x_G\), preserving its support, so every summand remains in \(J\). The atom-linear relations are closed. The odd derivation has square zero on generators and hence everywhere. This is a general DGA proof, not merely a calculation in rank two.\par
Write \(z=e_t\) in the affine model \(\widehat B\), where \(dz=x_t\). Adding top never changes whether a support is nested, so as graded algebras \(\widehat B=B_{\mathrm{pr}}\otimes\Lambda(z)\). The inclusion of \(B_{\mathrm{pr}}\) is a chain map; the graded projection \(q(z)=0\) generally is not.\par
Equation reference: eq:chain. For an atom \(H\), put \[\eta_H=\sum_{\substack{G\geq H\\G\ne t}}e_G,\qquad
 d\eta_H=-x_t.\] Let \(\phi_{\mathrm{aff}}\) be the affine algebra map used by the source, \(e_H\mapsto\sum_{G\geq H}e_G\), and let \(\phi_{\mathrm{pr}}=q\phi_{\mathrm{aff}}\) as a graded map. The derivation rule gives \[
 d\phi_{\mathrm{pr}}(\alpha)
   =-x_t\,\phi_{\mathrm{pr}}(\delta\alpha).\] Consequently its restriction to \(\overline{\mathrm{OS}}=\ker\delta\) is a chain map. An individual affine generator has differential \(-x_t\); it is not a projective element. Requiring every \(\eta_H\) to be closed would test the wrong domain.\par
There is also a precise conditional reduction of the quasi-isomorphism. Choose an atom \(H_0\), set \(v=z+\eta_{H_0}\), and observe \(dv=0\). Then \[\widehat B\simeq B_{\mathrm{pr}}\otimes\Lambda(v),\qquad
 \mathrm{OS}\simeq\overline{\mathrm{OS}}\otimes\Lambda(e_{H_0}).\] The second decomposition follows from \(\delta(e_{H_0})=1\): write \(\alpha=(\alpha-e_{H_0}\delta\alpha)+e_{H_0}\delta\alpha\). Under these decompositions, \(\phi_{\mathrm{aff}}=\phi_{\mathrm{pr}}\otimes\mathrm{id}\). Over \(\mathbb Q\), the external affine quasi-isomorphism therefore implies the corrected projective objectwise quasi-isomorphism. In rank one the proper exterior set is empty and the atom relation forces \(x_t=0\); the formula gives \(\mathbb Q\). The rank-zero tensor unit is treated separately.\par
In the \(U_{2,3}\) example, \(B_{\mathrm{pr}}\) has basis \(1,e_a,e_b,e_c,h\), and its only nonzero differential is \((-1,-1,-1):\mathbb Q^3\to\mathbb Qh\). Thus its cohomology dimensions are \(1,2,0\), as required for the projective OS algebra. Adding \(x_t=0\) instead would give three degree-one classes; it is not an adequate repair.\par
The local bar differential agrees with this correction: contracting a proper node inserts \(x_G\), as in \(\mathrm{FY}^{PD}([G])\) on p. 39, and moving that exterior generator to the contracted position gives the usual alternating sign. There is no projective top contraction. The underlying graded quotient is unchanged by [eq:bpr] (\#eq:bpr), so the cited nested-set decomposition has the same indexing.\par
These checks fix the DGA and its objectwise projective map. They do not independently prove the external affine Leray theorem, the complete graded decomposition, every odd-cooperad compatibility, or the identification with the canonical quadratic-dual map. The atom choice above is not claimed to give a natural cooperadic splitting. Corollary 6.6 needs a synchronized corrected comparison. Its conclusion is not disproved, and Corollary 6.10 has a separate Gröbner route.\par
\section*{Major F5b: the all-object projective OS assignment}
On p. 58, the source assigns \(\overline{\mathrm{OS}}(L)[-1]\) to every built lattice, while the construction is supposed to be a strong monoidal cooperad. Let each of two irreducible rank-one factors be \(B_1\). Each assigned vector space is one-dimensional in shifted degree one, so their tensor is one-dimensional in degree two. The product lattice \(B_2\) has projective OS dimensions \(1,1\) in original degrees \(0,1\), hence dimensions \(1,1\) in shifted degrees \(1,2\). Its assigned total dimension is two, not one. The rank-zero unit likewise receives a shifted copy of the ground field instead of the tensor unit.\par
No choice of structural maps can repair this object-level dimension obstruction. Define the displayed value only on irreducible arities, extend to a product by tensoring the factor values, and send the empty product to the tensor unit. This is a complete repair of the object assignment. It leaves the separate braiding issue below to be checked. The classification is Major because the literal core cooperad definition cannot be strong monoidal, not because the intended aritywise construction is refuted.\par
\section*{Five local issue groups}
\subsection*{F5a: projective FY presentations require irreducibility}
Proposition 2.23 and the following top-degree assertion (pp. 14–15) are stated for every built lattice. For \(B_2\) with reducible building set \(\{a,b\}\), Definition 2.22 gives \[\mathrm{FY}=\mathbb Z[x_a,x_b]/(x_a,x_b)=\mathbb Z.\] The displayed projective presentation instead gives \(\mathbb Z[x_a,x_b]/(x_a-x_b)\simeq\mathbb Z[u]\), since all its subsets are nested. There is no top generator in this building set.\par
Restrict the projective elimination and its top generator to irreducible built lattices. For a product with \(m\) maximal building factors the top degree is \(2(\operatorname{rk}L-m)\), by tensoring the irreducible factor rings. The later generator theorem explicitly uses irreducible arities. This is a Minor scope correction.\par
\subsection*{F6: use the least node greater than or equal to the input}
The general \(\mathrm{FY}(S)\) formula on p. 32 selects \(G'=\min S_{>G}\). It is undefined for \(G=t\). For maximal \(B_3\), \(S=\{ab,t\}\), \(G=ab\), it also selects the wrong interval: its alleged index is the bottom \(ab\) of \([ab,t]\), whereas the elementary map puts \(h_{ab}\) in \([0,ab]\).\par
Replace \(>\) by \(\geq\). Spanning gives existence and the nested-factor property gives a unique least node. If \(G\leq\tau_S(G')\), the building-factor decomposition would place \(G\) below a smaller node, contradicting minimality. Thus \(\tau_S(G')\vee G\) is a valid non-bottom index. Equality \(G'=G\), including top, now works. The elementary maps on p. 31 and their relations independently define the cooperad, so this is a Minor closed-form correction.\par
\subsection*{F7: only equality of fiber minima is needed}
In Lemma 5.25 (p. 54), the second claimed fiber equality is too strong. Take \(U_{2,3}\oplus U_{1,1}\) on atoms \(a<b<c<d\), with maximal building set, and the increasing chain \(S=\{a,abc,t\}\). Put \(G_1=a\), \(G_2=abc\). Then \(H_{G_2}=b\), \(\tau_{S\setminus\{G_1\}}(G_2)=0\), and \(\tau_S(G_2)=a\). The two displayed fibers are respectively \(\{b\}\) and \(\{b,c\}\), not equal.\par
For the required repair, put \(u=\tau_{S\setminus\{G_1\}}(G_2)\), \(v=\tau_S(G_2)\), \(h=H_{G_2}\). The fiber \(u\vee H=u\vee h\) is contained in the fiber \(v\vee H=G_2\), and contains \(h\). Since \(h\) is the minimum of the latter fiber, their minima agree. That is all inequality (24) uses. The first fiber comparison follows from the nested product decomposition. The increasing-label conclusion survives; this is a Minor unnecessary auxiliary equality.\par
\subsection*{F8: two outlook axioms need qualifications}
Definition 7.1 (p. 66) omits rank monotonicity. On \(E=\{a,b\}\), set \(r(\varnothing)=0\), \(r(a)=r(b)=1\), \(r(E)=0\). Every printed submodular and one-element upper-increment inequality holds, but \(r(a)>r(E)\). Add monotonicity, equivalently the lower bound in \(0\leq r(A\cup\{x\})-r(A)\leq1\).\par
The equivalence following Definition 7.2 is also not injective as stated. On the genuine simple loopless matroid \(U_{2,3}\), let \(\nu\) be discrete on proper subsets and one-block on \(E\); let \(\nu'\) be one-block on every nonempty subset. Both satisfy refinement, the block idempotence condition and rank additivity. They agree on all flats, namely the empty set, singletons and \(E\). They give the same building set but differ on pairs.\par
Equation reference: eq:closure. For the asserted bijection, impose \[
 \nu(X)=\nu(\operatorname{cl}X)|_X.\] On a flat, rank-additive blocks give a direct-sum restriction, so the blocks are closed. Refinement shows that every one-block flat below it lies in one such block; those are exactly its maximal building factors. Conversely building factors partition each flat. Their intersections with a spanning subset span their respective factors, so the induced partitions satisfy block idempotence and rank additivity. Equation [eq:closure] (\#eq:closure) determines every nonflat subset. This verifies the local bijection repair in the simple loopless setting. Alternatively the original decompositions can be described as richer data without claiming equivalence. Neither repair is a proof of the whole proposed future category. These outlook issues form one Minor group and are not used in the preceding geometric-lattice results.\par
\subsection*{F9: the lower polymatroid branch has a wrong index}
On p. 68, the lower branch sends \(h_{G'}\) to \(1\otimes h_G\) when \(G'\leq G\). In maximal \(B_3\), take \(G=ab\), \(G'=a\). The source has \(h_a=0\), while the printed image is \(1\otimes h_{ab}\ne0\) in the lower rank-two FY ring \(\mathbb Q[u]/(u^2)\). Thus the displayed algebra map is not well-defined even in the matroid special case. Replace the lower index by \(G'\), as in the earlier elementary FY formula. This is a Minor index correction; the whole polymatroid extension is not thereby certified.\par
\section*{A qualified shifted-sign concern}
This calculation uses the standard Koszul symmetry on graded vector spaces, consistent with the explicit graded-commutativity convention on p. 16. For \(B_3\), \(\mathcal G=\{a,b,c,t\}\), let \[\alpha=\delta(e_c e_b e_a)
       =e_b e_a-e_c e_a+e_c e_b\in\overline{\mathrm{OS}}(B_3).\] It is nonzero. Apply the displayed p. 58 maps literally. Cutting at \(a\) gives \((e_{ab}-e_{ac})\otimes1\); then cutting the upper factor at \(ab\) gives \(-1\otimes1\otimes1\). Cutting first at \(b\) gives \((e_{bc}-e_{ab})\otimes1\); the next cut at \(ab\) gives \(+1\otimes1\otimes1\). The final \(\sigma_{2,3}\) exchanges two shifted units of degree one, hence supplies another minus sign. The aligned composites are therefore equal and nonzero, where relation (27) requires them to be negatives.\par
An unsigned flip instead makes this particular relation hold. The source does not separately specify such a target symmetry in this construction. This report does not claim that changing the flip supplies a globally coherent repair, nor that every possible intended suspension convention fails. The exact calculation is a conditional failure of the displayed maps under the stated conventional interpretation. A fully specified symmetry/twist and recheck of all odd relations are needed. It is recorded as Needs Evidence for the intended convention, not counted as an additional convention-independent Major defect.\par
\section*{Other proof checks and the two Koszul routes}
The following positive checks delimit the findings. Building-factor decompositions justify induced intervals, the unique parent of a nonzero building element, and the maximality and injectivity of \(\operatorname{Comp}\). Nestedness ensures every structural local interval is nontrivial. Rank zero is the separately adjoined monoidal unit, not an irreducible structural arity. The proper grading can be checked directly as the number of source irreducible factors minus the number of target irreducible factors; it is additive under composition and tensor products.\par
The elementary FY maps preserve atom relations by semimodularity: an atom outside \(G\) joins \(G\) to an atom of the upper interval. Images of antichain relations reduce to the maximal image antichain, with additional factors multiplying an existing relation. Comparable cuts satisfy the three generator cases; nested incomparable cuts use the product-factor decomposition. The affine OS maps similarly kill circuit boundaries because either the lower subset is dependent or its upper contracted images are dependent. Affine OS is monoidal on product lattices, which does not imply the incorrect all-object projective assignment.\par
For Proposition 4.4, the coefficient-isolation induction has diagonal coefficient one. Read as a unitriangular integral argument, it gives a split injection on dual data and hence the required surjectivity. Mere injectivity of a dual map over \(\mathbb Z\) would not suffice; no torsion counterexample is inferred when this stronger argument is available. The degree identity \[\sum_{G\in S}\operatorname{rk}[\tau_S(G),G]=\operatorname{rk}L\] follows by peeling a minimal node and using the interval product. It checks the degree bounds and top-minus-one coefficient computations.\par
The shuffle construction uses the least atom in each join fiber, compatible with induced intervals. The geometric-lattice EL-chain property follows by greedily taking the least atom outside the current flat: delaying a smaller available atom forces a later decrease in labels. With corrected frames and fiber minima, this supplies the intended normal-monomial count in the Gröbner proof. The source’s reverse-prefix word order must be retained; substituting ordinary prefix lexicographic order changes its convention.\par
For Proposition 6.9 (p. 65), fix a decorated nested monomial. Quadraticity makes admissible contractions a condition on its adjacent edges. Normal-block decompositions form the augmented simplex on the admissible-edge set. If that set is nonempty, adjoining one fixed vertex gives a contracting homotopy. If it is empty, one fully split class remains on the quadratic diagonal. The filtration is bounded on each fixed finite arity/weight component relevant here. Thus this combinatorial mechanism supports Corollary 6.10 after the category and frame repairs. It does not use the ill-defined quotient differential of Definition 6.4.\par
The alternate Corollary 6.6 route needs the corrected DGA, the external Leray theorem, and synchronized cooperadic and suspension conventions. The two routes must not be conflated. Full proofs of the cited FY basis theorem, the Bibby–Denham–Feichtner Leray comparison, the Kaufmann–Ward/Ward bar–cobar framework, the Hodge results and geometric identifications were not independently retrieved or reproved. Their use and the local hypotheses were tracked in the review evidence. The general nonhomogeneous or infinite-dual extensions are not certified beyond the checked finite FY components. Unavailable code is not an error.\par
\section*{Calculations, repair status, and assessment}
The lead reviewer’s exact Python calculation verified the building axiom on all eight Boolean elements, all 29 node orderings for the two stated \(B_3\) building sets, the six-versus-three frame count, the nonzero Leray top class and differential rank, the two-element rank inequalities, both \(U_{2,3}\) partition systems, the unequal fibers, and the conditional OS signs. Exact augmented simplex boundary matrices for zero through five vertices satisfy \(d^2=0\) and have the stated homology, including the empty case. The process exited with status zero. These are finite checks supporting explicit examples; the general repairs are the arguments above, not extrapolations from enumeration.\par
The normal-form helper additionally checked 19 local refinements, 64 pointwise \(\operatorname{Comp}\)–\(\Phi\) identities and 186 naturality instances. The independent Leray helper supplied the general support-ideal proof and the conditional affine/projective splitting. All final mathematical evidence is tied to the exact source and reviewer files; author programs were not executed.\par
Source item | Assessment | Verified repair | Remaining scope\par
F1, Lemma 3.20 | Major | Building-set invariant and full recovery | No false normal-form conclusion shown\par
F2, p. 53 frames | Major | Allow rank-one top interval | Corrected basis argument\par
F3, p. 57 action | Major | Insert in current parent interval | Corrected category required\par
F4, Def. 6.4 | Major | Proper exterior generators; projective chain map | External/global comparison\par
F5b, p. 58 objects | Major | Irreducible values, tensor extension | Odd symmetry still to specify\par
F5a, F6–F9 | Minor | Scope, index, minima and outlook corrections above | Future extensions not certified\par
Shifted signs | Needs Evidence | Conditional calculation only | No global sign repair claimed\par
There are five demonstrated Major issue groups and five Minor groups; F5a and F5b separate two statements with different roles. The conditional sign observation is additional and is not counted as a demonstrated convention-independent error. No example here disproves the intended final Koszulness conclusion. Author adoption of all repairs is unknown. The scientific decision remains pending; human reviews are empty.\par
9 B. Coron, Matroids, Feynman categories, and Koszul duality, arXiv:2211.12370v3, 11 December 2024, 70 pp. Exact supplied PDF identified in Section 1. Assigned journal DOI: 10.1215/00127094-2024-0077.\par
\end{document}
